% GENERATED FILE. Edit canonical lessons, then run npm run generate:books.
% canonical-source-hash: fnv1a64:43c1370ba0879056
% canonical-lessons: HI-C498-parvaris, HI-C498-lori, HI-C498-pilana, HI-C498-khilana, HI-C498-palnaintr, HI-R498-review-492-498

\chapter{Bringing Up a Child}
\label{ch:hi-bringing-up-a-child}
\hlchaptermodality{498}

\begin{chapteropening}
\textbf{By the end of this chapter:} \emph{I can say how a child was brought up.}

Complete the last lesson of chapter 498: I can say how a child was brought up.
\end{chapteropening}

\section[\texorpdfstring{\hi{परवरिश} (parvariś)}{परवरिश (parvariś)}]{\hi{परवरिश} (parvariś) — nurture; upbringing}

\label{lesson:HI-C498-parvaris}



\begin{quote}
Before the new one: say the Hindi for newborn, then the Hindi for to be born.
\end{quote}

\subsection*{You'll want to know: \hi{परवरिश}}
\textbf{\hi{परवरिश}} — \emph{parvariś} — \textquotedblleft{}nurture; upbringing\textquotedblright{}.

\begin{quote}
\hi{दादी} \hi{ने} \hi{अरुण} \hi{की} \hi{परवरिश} \hi{की।}
\end{quote}

\begin{quote}
\emph{dādī ne Aruṇ kī parvariś kī.}
\end{quote}

\begin{quote}
Grandmother brought Arun up.
\end{quote}

Feminine.

\begin{cousinweb}
\textbf{\hi{परवरिश}} (\emph{parvariś}) is usually traced to Persian \emph{parvarish}, \textquotedblleft{}nurture, fostering\textquotedblright{}.
\end{cousinweb}

\begin{grammarlens}[title={What you've built}]
One more word for describing what you lived and felt.
\end{grammarlens}

\subsection*{Your turn}
\begin{itemize}
\raggedright
  \item \emph{Say it:} \emph{parvariś}
  \item \emph{Say it:} \emph{parvariś}, once more
  \item \emph{Say it:} \emph{dādī ne Aruṇ kī parvariś kī.} — Grandmother brought Arun up.
  \item \emph{Recall:} say the Hindi for newborn, then the Hindi for to be born, then say \emph{parvariś} again
\end{itemize}

\begin{tcolorbox}[breakable,colback=teal!4,colframe=teal!35!black,title={Before you move on}]
What does \hi{परवरिश} mean? (\textquotedblleft{}Nurture; upbringing\textquotedblright{}.) Say it once more.
\end{tcolorbox}
\section[\texorpdfstring{\hi{लोरी} (lorī)}{लोरी (lorī)}]{\hi{लोरी} (lorī) — a lullaby}

\label{lesson:HI-C498-lori}



\begin{quote}
Before the new one: say the Hindi for to be born, then the Hindi for nurture.
\end{quote}

\subsection*{You'll want to know: \hi{लोरी}}
\textbf{\hi{लोरी}} — \emph{lorī} — \textquotedblleft{}a lullaby\textquotedblright{}.

\begin{quote}
\hi{माँ} \hi{ने} \hi{लोरी} \hi{गाकर} \hi{बच्चे} \hi{को} \hi{सुलाया।}
\end{quote}

\begin{quote}
\emph{māṁ ne lorī gākar bacce ko sulāyā.}
\end{quote}

\begin{quote}
Mother sang a lullaby and got the baby to sleep.
\end{quote}

Feminine.

\begin{cousinweb}
\textbf{\hi{लोरी}} (\emph{lorī}) is a tadbhava: CDIAL lists Hindi \emph{lorī}, \textquotedblleft{}lullaby\textquotedblright{}, under \emph{*lōḍa-}, a form it reconstructs (CDIAL 11136).
\end{cousinweb}

\begin{grammarlens}[title={What you've built}]
One more word for describing what you lived and felt.
\end{grammarlens}

\subsection*{Your turn}
\begin{itemize}
\raggedright
  \item \emph{Say it:} \emph{lorī}
  \item \emph{Say it:} \emph{lorī}, once more
  \item \emph{Say it:} \emph{māṁ ne lorī gākar bacce ko sulāyā.} — Mother sang a lullaby and got the baby to sleep.
  \item \emph{Recall:} say the Hindi for to be born, then the Hindi for nurture, then say \emph{lorī} again
\end{itemize}

\begin{tcolorbox}[breakable,colback=teal!4,colframe=teal!35!black,title={Before you move on}]
What does \hi{लोरी} mean? (\textquotedblleft{}A lullaby\textquotedblright{}.) Say it once more.
\end{tcolorbox}
\section[\texorpdfstring{\hi{पिलाना} (pilānā)}{पिलाना (pilānā)}]{\hi{पिलाना} (pilānā) — to give someone a drink}

\label{lesson:HI-C498-pilana}



\begin{quote}
Before the new one: say the Hindi for nurture, then the Hindi for a lullaby.
\end{quote}

\subsection*{You'll want to know: \hi{पिलाना}}
\textbf{\hi{पिलाना}} — \emph{pilānā} — \textquotedblleft{}to give someone a drink\textquotedblright{}.

\begin{quote}
\hi{माँ} \hi{ने} \hi{बच्चे} \hi{को} \hi{दूध} \hi{पिलाया।}
\end{quote}

\begin{quote}
\emph{māṁ ne bacce ko dūdh pilāyā.}
\end{quote}

\begin{quote}
The mother gave the baby some milk.
\end{quote}

Past, with \hi{ने} (he or she …): \hi{उसने} \hi{बच्चे} \hi{को} \hi{पानी} \hi{पिलाया}, \emph{usne bacce ko pānī pilāyā}.

\begin{cousinweb}
\textbf{\hi{पिलाना}} (\emph{pilānā}) is the causative partner of \emph{pīnā}, \textquotedblleft{}to drink\textquotedblright{}, which CDIAL lists under Sanskrit \emph{píbati}, \textquotedblleft{}drinks\textquotedblright{} (CDIAL 8209): to make someone drink.
\end{cousinweb}

\begin{grammarlens}[title={What you've built}]
One more word for describing what you lived and felt.
\end{grammarlens}

\subsection*{Your turn}
\begin{itemize}
\raggedright
  \item \emph{Say it:} \emph{pilānā}
  \item \emph{Say it:} \emph{pilānā}, once more
  \item \emph{Say it:} \emph{māṁ ne bacce ko dūdh pilāyā.} — The mother gave the baby some milk.
  \item \emph{Recall:} say the Hindi for nurture, then the Hindi for a lullaby, then say \emph{pilānā} again
\end{itemize}

\begin{tcolorbox}[breakable,colback=teal!4,colframe=teal!35!black,title={Before you move on}]
What does \hi{पिलाना} mean? (\textquotedblleft{}To give someone a drink\textquotedblright{}.) Say it once more.
\end{tcolorbox}
\section[\texorpdfstring{\hi{खिलाना} (khilānā)}{खिलाना (khilānā)}]{\hi{खिलाना} (khilānā) — to feed}

\label{lesson:HI-C498-khilana}



\begin{quote}
Before the new one: say the Hindi for a lullaby, then the Hindi for to give someone a drink.
\end{quote}

\subsection*{You'll want to know: \hi{खिलाना}}
\textbf{\hi{खिलाना}} — \emph{khilānā} — \textquotedblleft{}to feed\textquotedblright{}.

\begin{quote}
\hi{दादी} \hi{ने} \hi{अपने} \hi{हाथ} \hi{से} \hi{बच्चों} \hi{को} \hi{खाना} \hi{खिलाया।}
\end{quote}

\begin{quote}
\emph{dādī ne apne hāth se baccoṁ ko khānā khilāyā.}
\end{quote}

\begin{quote}
Grandmother fed the children with her own hands.
\end{quote}

Past, with \hi{ने} (he or she …): \hi{उसने} \hi{बच्चे} \hi{को} \hi{खाना} \hi{खिलाया}, \emph{usne bacce ko khānā khilāyā}.

\begin{cousinweb}
\textbf{\hi{खिलाना}} (\emph{khilānā}) is the causative partner of \emph{khānā}, \textquotedblleft{}to eat\textquotedblright{}, which CDIAL lists under Sanskrit \emph{khādati}, \textquotedblleft{}eats\textquotedblright{} (CDIAL 3865): to make someone eat. Another \emph{khilānā}, \textquotedblleft{}to make bloom\textquotedblright{}, belongs to \emph{khilnā}.
\end{cousinweb}

\begin{grammarlens}[title={What you've built}]
One more word for describing what you lived and felt.
\end{grammarlens}

\subsection*{Your turn}
\begin{itemize}
\raggedright
  \item \emph{Say it:} \emph{khilānā}
  \item \emph{Say it:} \emph{khilānā}, once more
  \item \emph{Say it:} \emph{dādī ne apne hāth se baccoṁ ko khānā khilāyā.} — Grandmother fed the children with her own hands.
  \item \emph{Recall:} say the Hindi for a lullaby, then the Hindi for to give someone a drink, then say \emph{khilānā} again
\end{itemize}

\begin{tcolorbox}[breakable,colback=teal!4,colframe=teal!35!black,title={Before you move on}]
What does \hi{खिलाना} mean? (\textquotedblleft{}To feed\textquotedblright{}.) Say it once more.
\end{tcolorbox}
\section[\texorpdfstring{\hi{पलना} (palnā)}{पलना (palnā)}]{\hi{पलना} (palnā) — to be brought up}

\label{lesson:HI-C498-palnaintr}



\begin{quote}
Before the new one: say the Hindi for to give someone a drink, then the Hindi for to feed.
\end{quote}

\subsection*{You'll want to know: \hi{पलना}}
\textbf{\hi{पलना}} — \emph{palnā} — \textquotedblleft{}to be brought up\textquotedblright{}.

\begin{quote}
\hi{अरुण} \hi{अपने} \hi{नाना} \hi{के} \hi{घर} \hi{पला।}
\end{quote}

\begin{quote}
\emph{Aruṇ apne nānā ke ghar palā.}
\end{quote}

\begin{quote}
Arun was brought up at his grandfather's house.
\end{quote}

Past (he …): \hi{वह} \hi{गाँव} \hi{में} \hi{पला}, \emph{vah gā̃v meṁ palā}.

\begin{cousinweb}
\textbf{\hi{पलना}} (\emph{palnā}) is a tadbhava: CDIAL lists Hindi \emph{palnā}, \textquotedblleft{}to be reared\textquotedblright{}, beside \emph{pālnā}, \textquotedblleft{}to rear\textquotedblright{}, under Sanskrit \emph{pāláyati}, \textquotedblleft{}protects\textquotedblright{} (CDIAL 8129). It is the intransitive partner of \emph{pālnā}.
\end{cousinweb}

\begin{grammarlens}[title={What you've built}]
One more word for describing what you lived and felt.
\end{grammarlens}

\subsection*{Your turn}
\begin{itemize}
\raggedright
  \item \emph{Say it:} \emph{palnā}
  \item \emph{Say it:} \emph{palnā}, once more
  \item \emph{Say it:} \emph{Aruṇ apne nānā ke ghar palā.} — Arun was brought up at his grandfather's house.
  \item \emph{Recall:} say the Hindi for to give someone a drink, then the Hindi for to feed, then say \emph{palnā} again
\end{itemize}

\begin{tcolorbox}[breakable,colback=teal!4,colframe=teal!35!black,title={Before you move on}]
What does \hi{पलना} mean? (\textquotedblleft{}To be brought up\textquotedblright{}.) Say it once more.
\end{tcolorbox}
\section[(dialogue)]{Looking back: from Getting Better to Bringing Up a Child}

\label{lesson:HI-R498-review-492-498}



\begin{quote}
Say the Hindi for a massage and healthy.
\end{quote}

\subsection*{Your turn}
Say these aloud:

\begin{itemize}
\raggedright
  \item a massage, healthy, well again, ayurveda, vigour
  \item diabetes, asthma, paralysis, cancer, to stiffen
  \item nursing the sick, a lap, to lay someone down, to warm (a sore part), to cover someone (with a blanket)
  \item exercise, yoga, avoiding certain foods, nourishing, to be digested
  \item insomnia, drowsy, to sprawl, to mumble, a bedroll
  \item birth, pregnant, childbirth, newborn, to be born
  \item nurture, a lullaby, to give someone a drink, to feed, to be brought up
\end{itemize}

\begin{tcolorbox}[breakable,colback=teal!4,colframe=teal!35!black,title={Before you move on}]
A massage? (\textbf{\hi{मालिश}}, \emph{māliś}.) Healthy? (\textbf{\hi{स्वस्थ}}, \emph{svasth}.) Well again? (\textbf{\hi{चंगा}}, \emph{caṅgā}.) Ayurveda? (\textbf{\hi{आयुर्वेद}}, \emph{āyurved}.) Vigour? (\textbf{\hi{स्फूर्ति}}, \emph{sphūrti}.) Diabetes? (\textbf{\hi{मधुमेह}}, \emph{madhumeh}.) Asthma? (\textbf{\hi{दमा}}, \emph{damā}.) Paralysis? (\textbf{\hi{लकवा}}, \emph{lakvā}.) Cancer? (\textbf{\hi{कैंसर}}, \emph{kaiṁsar}.) To stiffen? (\textbf{\hi{जकड़ना}}, \emph{jakaṛnā}.) Nursing the sick? (\textbf{\hi{तीमारदारी}}, \emph{tīmārdārī}.) A lap? (\textbf{\hi{गोद}}, \emph{god}.) To lay someone down? (\textbf{\hi{लिटाना}}, \emph{liṭānā}.) To warm (a sore part)? (\textbf{\hi{सेंकना}}, \emph{sẽknā}.) To cover someone (with a blanket)? (\textbf{\hi{ओढ़ाना}}, \emph{oṛhānā}.) Exercise? (\textbf{\hi{कसरत}}, \emph{kasrat}.) Yoga? (\textbf{\hi{योग}}, \emph{yog}.) Avoiding certain foods? (\textbf{\hi{परहेज़}}, \emph{parhez}.) Nourishing? (\textbf{\hi{पौष्टिक}}, \emph{pauṣṭik}.) To be digested? (\textbf{\hi{पचना}}, \emph{pacnā}.) Insomnia? (\textbf{\hi{अनिद्रा}}, \emph{anidrā}.) Drowsy? (\textbf{\hi{उनींदा}}, \emph{unīndā}.) To sprawl? (\textbf{\hi{पसरना}}, \emph{pasarnā}.) To mumble? (\textbf{\hi{बुदबुदाना}}, \emph{budbudānā}.) A bedroll? (\textbf{\hi{बिछौना}}, \emph{bichaunā}.) Birth? (\textbf{\hi{जन्म}}, \emph{janm}.) Pregnant? (\textbf{\hi{गर्भवती}}, \emph{garbhvatī}.) Childbirth? (\textbf{\hi{प्रसव}}, \emph{prasav}.) Newborn? (\textbf{\hi{नवजात}}, \emph{navjāt}.) To be born? (\textbf{\hi{पैदा} \hi{होना}}, \emph{paidā honā}.) Nurture? (\textbf{\hi{परवरिश}}, \emph{parvariś}.) A lullaby? (\textbf{\hi{लोरी}}, \emph{lorī}.) To give someone a drink? (\textbf{\hi{पिलाना}}, \emph{pilānā}.) To feed? (\textbf{\hi{खिलाना}}, \emph{khilānā}.) To be brought up? (\textbf{\hi{पलना}}, \emph{palnā}.)
\end{tcolorbox}
